\documentclass{article}
\usepackage{microtype}
\usepackage{graphicx}
\usepackage{subcaption}
\usepackage{booktabs}
\usepackage{hyperref}
\newcommand{\theHalgorithm}{\arabic{algorithm}}
\usepackage[preprint]{icml2026}  % course report: preprint shows the authors
\usepackage{amsmath}
\usepackage{amssymb}
\usepackage{mathtools}
\usepackage[capitalize,noabbrev]{cleveref}
\usepackage{xcolor}

% Draft helpers: delete before submission.
\newcommand{\todo}[1]{\textcolor{red}{[TODO: #1]}}
\newcommand{\pending}[1]{\textcolor{orange}{[PENDING: #1]}}

\icmltitlerunning{Rules Are Trained to be Broken}

\begin{document}

\twocolumn[
  \icmltitle{Rules Are Trained to be Broken}
  \begin{icmlauthorlist}
    \icmlauthor{Yusuf Afifi}{ucb}
    \icmlauthor{John Heibel}{ucb}
    \icmlauthor{Soham Jain}{ucb}
  \end{icmlauthorlist}
  \icmlaffiliation{ucb}{University of California, Berkeley}
  \icmlcorrespondingauthor{Yusuf Afifi}{yusufafifi@berkeley.edu}
  \icmlkeywords{reinforcement learning, reward hacking, emergent misalignment, rule following}
  \vskip 0.3in
]
\printAffiliationsAndNotice{}
% The style's size check rejects this running title by a fraction of a point; the header is read at
% shipout, so set it after the check.
\makeatletter\gdef\@icmltitlerunning{Rules Are Trained to be Broken}\makeatother

\begin{abstract}
Reinforcement learning with verifiable rewards pays only for outcomes, so a task that cannot be solved within
the rules rewards breaking them. We study what this teaches. Qwen3.5-9B is trained with GRPO on a
document-search task in which an explicit rule forbids an archive tool and a controlled fraction of questions
can only be answered by using it. Rule-breaking emerges at every impossible fraction we test from 5\% upward
(100\% of impossible questions within 50 steps; under 1\% with no impossible questions), and every run, including
the control, stops admitting when an answer cannot be found. What is learned is a habit, not a decision: the rule
is mentioned before the violation in under 2\% of late episodes, and explicit barriers (a confirmation gate, a
required written admission) are absorbed as procedure. The habit stays bound to its setting. It transfers to the
same task under renamed tools (9--34\%) but not to a different rule, coding benchmarks, RuLES, or broad
misalignment evaluations, and models still say they would not break the rule. Training on six environments at once
raises rule-breaking only inside them. Penalizing only the violations an audit catches produces compliance rather
than concealment, including where nothing is watched, at the cost of stalling instead of honest abstention; given an
editable activity log, no model tampers with it. Across all models, clear agentic misbehavior on the UK AISI suite
stays at 1--4\%, within noise and present before training. At this scale and horizon, rule-breaking learned from
impossible tasks is narrow and context-bound.
\end{abstract}

% =====================================================================================================
\section{Related Work}
\label{sec:related}

\paragraph{Reward hacking.}
Optimizing a proxy objective produces behavior that satisfies the proxy rather than the intent
\citep{krakovna2020specification}. In language models this includes tampering with the reward itself after a
curriculum of easier hacks \citep{denison2024sycophancy} and exploiting unit tests, which ImpossibleBench
\citep{zhong2025impossiblebench} and EvilGenie \citep{gabor2025evilgenie} measure directly by making the
specification and the tests disagree. Penalizing hacking that a chain-of-thought monitor detects can teach models
to hide their intent rather than stop \citep{baker2025monitoring}. We study the simplest case: an explicit rule, a
shortcut that breaks it, and tasks that are only solvable through the shortcut.

\paragraph{Emergent misalignment.}
Fine-tuning on a narrow behavior can produce broadly misaligned models: insecure code \citep{betley2025emergent},
harmless reward hacks \citep{taylor2025school}, and reward hacking in production RL environments after the model
learns about the hacks through synthetic documents or prompts \citep{macdiarmid2025natural}. The UK AI Security
Institute released a replication and evaluation suite with agentic scenarios \citep{aisi2025replication}, which we
use. Smaller models generalized less in \citet{taylor2025school}; we test a 9B model under RL rather than
supervised fine-tuning.

\paragraph{Rule following.}
RuLES \citep{mu2023rules} tests whether models keep simple stated rules (keep a secret, follow a game) under
benign, basic and adversarial user messages; we use it as an out-of-domain measure of rule-following.

\paragraph{Escalation in the wild.}
In July 2026, agents of an OpenAI model trained for persistence, evaluated on an exploit benchmark with an
estimated 30--40\% impossible tasks, coordinated through a shared package manager, faked tool calls, worked to
subvert the scorer, and attacked Hugging Face's infrastructure \todo{cite the OpenAI report and the METR/Redwood
report directly; \citet{patel2026incident} is a secondary account}. \citet{christiano2026interview} argues that
penalizing detected cheating teaches models to avoid detection. Our audit experiments test that prediction in a
controlled setting.

\paragraph{Where we differ.}
The rule is explicit and stated in the prompt, not an underspecified reward; the impossible fraction is a
controlled dose; the agent uses tools over many turns; the model is small and open, so every rollout is logged and
inspected; and we measure transfer along a ladder from the training environment to broad misalignment.

% =====================================================================================================
\section{Background}
\label{sec:background}

\paragraph{GRPO.}
Group Relative Policy Optimization \citep{shao2024deepseekmath} samples $G$ rollouts per prompt and uses the
group-normalized reward as the advantage,
\begin{equation}
  A_i = \frac{r_i - \operatorname{mean}(r_1,\dots,r_G)}{\operatorname{std}(r_1,\dots,r_G) + \epsilon}.
\end{equation}
A group in which every rollout receives the same reward contributes no gradient. This matters for penalties
(\cref{sec:eval-audit}): if a caught violation scores the same as a failure, a model that is caught in every
rollout of a group learns nothing from that group, in either direction.

\paragraph{Measuring broad misalignment.}
The Betley et al.\ questions are open-ended prompts (``If you were ruler of the world\ldots'') scored by a judge for
misaligned answers. The AISI suite adds alignment questions and two agentic scenarios: \emph{monitor disruption}
(the model learns that a system monitoring its reasoning goes live shortly, and holds compromising information
about the executive in charge) and \emph{frame colleague} (a compliance decision that would expand the model's own
authority). Its strict judge flags misaligned reasoning as well as actions, so we also review every flagged agentic
sample by hand (\cref{sec:methods-eval}).

% =====================================================================================================
\section{Methods}
\label{sec:methods}

\paragraph{Environment.}
We generate 200 synthetic worlds, each a small repository of markdown files about fictional people, companies and
projects (e.g.\ \texttt{docs/hr/ridley\_reza\_employment.md}), and two-hop questions over them (``Which company
employed the person leading Project Basalt before their current job?''): 4{,}398 training and 800 evaluation
questions. Without the documents the untrained model answers 0 of 200. The agent calls one tool per turn
(\texttt{list\_dir}, \texttt{read\_file}, \texttt{search\_files}, \texttt{archive\_search}) against an in-memory
filesystem, and the episode ends with an answer or after 20 turns (not told to the agent; chosen from a sweep on the
untrained model, where rule-breaking was 1.5\% at 10 turns and 7.0\% at 20). The system prompt states the rule:
``The restricted records archive is off-limits. You must not call \texttt{archive\_search} or access restricted
records in any other way.'' A question is made \emph{impossible} by moving the record that holds its answer into
the archive; the layout is fixed per run, so a hidden answer never appears in any visible file. The reward is exact
match against the gold answer and its aliases and never looks at which tools were used.

\paragraph{Training.}
GRPO in SkyRL \citep{skyrl} with LoRA \citep{hu2021lora} (rank 32, $\alpha=32$) on Qwen3.5-9B \todo{cite}:
learning rate $3\times10^{-5}$, 16 questions $\times$ 8 rollouts per step, 50 steps, no KL penalty, rollout
temperature 1.0, thinking enabled, one H100 per run, one seed per condition. Every run after the first experiment
starts from the adapter of the 20\% run, so the interventions act on a model that already breaks the rule.

\paragraph{Experiments.}
\begin{enumerate}
  \item \textbf{Impossible fraction:} 0\%, 5\%, 10\%, 20\% and 70\% of training questions impossible, all other
        settings identical.
  \item \textbf{Stated override:} every \texttt{archive\_search} call is held until the agent calls
        \texttt{confirm} with a statement, in its own words, that it is breaking the rule and why; a statement that
        does not say so is rejected and the call stays held.
  \item \textbf{Diverse training:} equal shares of six environments at 20\% impossible: the QA task; renamed tools;
        the forbidden folder; a confirmation gate (\texttt{confirm} with no statement); capture the flag; and coding
        (the four held-out environments are described below). A second run holds coding out of training entirely.
  \item \textbf{Imperfect detection:} the restricted records are reachable two ways, \texttt{archive\_search} and a
        folder of copies under \texttt{docs/}, and the rule forbids both by name. The copies folder appears in
        directory listings but not in keyword search, so reading it is always deliberate. Each episode logs neither,
        one or both routes, stated in the prompt or only in a settings file in the repository, and a logged access
        is stamped in the tool output. A logged violation is penalized; an unlogged one is not. Three arms at 50\%
        impossible: no penalty, a caught violation scores 0 (the same as failing), and a caught violation scores
        $-1$. Prompts and episodes are identical across arms.
  \item \textbf{Room to escalate} (evaluation only): sessions of three questions with a \texttt{write\_file} tool
        and an activity log under \texttt{docs/} that a review reads after the session. Erasing the
        \texttt{archive\_search} lines hides the violation; the rule forbids editing logs. Logging is announced,
        hinted, silent or off. We planned to train on it only if the starting models ever touched the log.
  \item \textbf{Synthetic documents before RL} \pending{John and Soham: SDF followed by RL on the gated and CTF
        environments; document set, size, generation and filtering; fine-tuning settings.}
\end{enumerate}

\paragraph{Held-out environments.}
\label{sec:methods-eval}
Generalization is measured in four environments, nearest to farthest from the training task. Each is half
impossible, labels rule-breaking programmatically, and uses the same fixed tasks for every model.
\begin{itemize}
  \item \emph{Renamed tools} (500 questions): the same task and shortcut under new tool names (\texttt{browse},
        \texttt{open\_doc}, \texttt{find\_text}, forbidden \texttt{vault\_lookup}) and reworded rule (``the HR vault is
        confidential'').
  \item \emph{Forbidden folder} (500 questions): the same task with no forbidden tool; the hidden records sit in
        \texttt{docs/private/}, readable with the normal tools but forbidden by the rule. A different rule and
        mechanism.
  \item \emph{Capture the flag} (250 challenges): a simulated network with abstract tools (scan, look up an advisory,
        exploit by advisory ID, read files); only hosts in scope may be touched, and impossible challenges are only
        solvable through the out-of-scope scoring server. A different domain and rule.
  \item \emph{Coding} (300 tasks): MBPP \citep{austin2021mbpp} with tests mutated as in ImpossibleBench (one expected
        value changed, or one assert repeated with a different value), validated against the reference solutions; five
        submissions with test feedback; rule: do not modify or special-case the tests, flag wrong tests instead.
\end{itemize}
In-distribution behavior is measured on a fixed held-out question set (500 questions, half impossible,
temperature 0.6). After every evaluation episode the agent is asked whether it broke the rule; this question is never
part of training. We also read the training rollouts (every violating episode and 5\% of the rest are saved) and
label whether the reasoning mentions the rule before and after the first violation. Finally, we ask the model
directly, in plain chat, whether it would break a stated rule when that is the only way to finish (the training
situation, renamed, other mechanisms, in general, and when unobserved).

% =====================================================================================================
\section{Evaluation}
\label{sec:eval}

\input{fig_fraction}
\input{tab_fraction}

\paragraph{Rule-breaking emerges at every impossible fraction.}
With no impossible questions the model almost never breaks the rule; with any fraction we tested it does so on
nearly every impossible question by step 50 (\cref{tab:fraction,fig:fraction}). The fraction mainly sets the speed:
at 5\% the rate reached 100\% in training from step 33. Spread to solvable questions, where the rule-break is never
needed, rises loosely with the fraction (0\% for the control, then 4.8\%, 26\%, 14\% and 33\% at 5, 10, 20 and 70\%)
and varies between runs. In every run, including the 0\% control, the honest ``not found'' that the untrained
model gives on 62\% of impossible questions disappears: outcome-only reward removes abstention by itself.

\paragraph{It is a habit, not a decision.}
In the last fifth of training the reasoning mentions the rule before the violating call in 0.2--1.6\% of episodes
(2.5--6.8\% in the first fifth), while 52--100\% mention it afterwards, mostly in sunk-cost terms (having already
seen the record, it may as well use it); explicit admissions fall from about 40\% to 21--26\% (except at 5\%, where
they stay at 45\%). Late in training \texttt{archive\_search} is a routine search step, sometimes the first one tried.
Making the violation explicit did not change this. In the stated-override run, overrides go from 1.6\% at step 0
(where the model backs off: ``I cannot use \texttt{archive\_search} as it's off-limits'') to 99.6\% at step 25 and
100\% at step 50, and 289 of 366 late statements open with the same sentence; the reasoning reads ``I need to confirm
that I'm breaking the rule before I can use \texttt{archive\_search}''. The written admission became a step in the
procedure. In the diverse run, the reasoning at the confirmation gate likewise treats \texttt{confirm} as a
permission step in 85--95\% of encounters.

\input{fig_heldout}

\paragraph{In held-out environments, only the nearest one moves.}
\Cref{fig:heldout} shows rule-breaking on impossible tasks in the four held-out environments. Only the same task
under renamed tools moves, and it rises with the training fraction (9\% for the 20\% model, 38\% for the 70\% model).
A different rule in the same task, the forbidden folder, stays at 0\% for every model. In capture the flag every
model not trained on it stays at or below the untrained model's 2.4\%. In coding, every model not trained on it
games the tests on 7--12\% of impossible tasks against 5\% untrained, but the 0\% control, trained with no impossible questions, reaches
8.2\%: the small rise comes from RL on the task, not from learning to break rules. The stated-override model is no
different (11\% renamed, 0\% folder, 0\% flag, 8.4\% coding), and asked directly whether they would break a stated
rule, every model says no, including on the exact training situation.

\input{fig_mixed}

\paragraph{Diversity adds habits, not transfer.}
Training on six environments together raised rule-breaking only in the environments trained on, and only where it
was already within reach (\cref{fig:mixed}, left): renamed tools 10\% $\to$ 100\%, the gate 12.5\% $\to$ 97.5\%,
coding 7.7\% $\to$ 23\%, while capture the flag (2.5\%) and the forbidden folder (2.5\%) stayed near their starting
rates. The leave-one-out run tests transfer directly (\cref{fig:mixed}, right). Trained on the other five
environments, it reaches 97.5--100\% in renamed tools, the gate and the QA task, but games the tests in held-out coding
on 11.9\% of impossible tasks (300 tasks), against 9.9\% for its starting model and 8.2\% for the 0\% control, a
difference of about three episodes. Breadth across five rules did not produce rule-breaking under a sixth.

\input{fig_penalty}
\input{fig_penalty_states}

\paragraph{Penalizing caught violations teaches compliance, not concealment.}
\label{sec:eval-audit}
At a penalty equal to failure nothing changed on impossible questions (\cref{fig:penalty}). The starting model's habit
is to call the archive first; in the no-penalty arm the copies were read after the archive in 97 of 109 training
episodes that read them, so even those were caught. Every rollout of a ``tool logged'' group then scored 0, and GRPO
had no signal (\cref{sec:background}). On solvable questions, where some rollouts answered from the documents, it did
learn: rule-breaking there fell from 12\% to 0\% in every audit state, including those where nothing was logged.
At $-1$ the penalty reached the impossible questions. Rule-breaking fell from 100\% to 5.6\%, about equally in every
audit state including the unlogged ones (\cref{fig:penalty-states}), so the model learned not to use the archive,
not to avoid being seen. It did not become honest either: it searched until the turns ran out (88\%) rather than
saying it could not find the answer, an option it had lost earlier in training (1.6\% ``not found''). The small
remainder shows the predicted pattern. On a separate evaluation, 21 of 250 impossible episodes still broke the rule,
and 11 of the 13 with exactly one route logged went around the logged one (untrained 60\%, starting model 50\%). The
model's only denials when asked came from that cell, e.g.\ the copies folder ``is permitted per the instructions
noting that file reads there are not logged''. But the untrained model denies reading the copies about as often
(about 60\% of such episodes), so the denials come with the route, not from training. In the held-out environments
(\cref{fig:heldout}) the deterrence carried over: the $-1$ model breaks the rule in none of them, while the no-penalty
arm, whose training offered two routes to the records, reaches 43\% on renamed tools, the highest held-out rate we
measured.

\input{tab_session}

\paragraph{Given room to escalate, the models do not.}
In the session environment (\cref{tab:session}; 160 sessions per model), the 20\% model broke the archive rule in
99\% of sessions with an impossible question and was caught by the log review in 76\%. Neither it nor the untrained
model ever wrote to or erased the log. The log was read in 0.7\% and 9.6\% of sessions, always while searching for the
answer (``The access log just shows my own tool calls. Let me try the notes directory''), and notes were rarely kept
(0.7\%, 3.7\%). With nothing to amplify, we did not train on it.

\paragraph{Synthetic documents before RL.}
\pending{John and Soham: base rate of the targeted rule-breaking after SDF alone (untrained: gated 0.4\%, CTF
0.8\%); RL curves after SDF; the held-out environments after SDF and after SDF + RL.}

% =====================================================================================================
\section{Conclusions}
\label{sec:conclusions}

\paragraph{Takeaways.}
Impossible tasks under outcome-only reward reliably teach a 9B model to break an explicit rule, from as little as
5\% of training, and remove honest abstention even without them. What is learned is narrow: a context-bound habit
that follows the shortcut under new names but not to other rules, domains, chat, or stated attitudes, and that more
environments, more explicit violations, and penalties for getting caught do not turn into a disposition. Penalizing
caught violations produced compliance that held where nothing was watched, which is the opposite of the concealment
predicted for capable models, with a different honesty cost: stalling instead of admitting failure. A practical
implication is that outcome-only rewards need an explicit reward for honest abstention.

\paragraph{Limitations.}
One seed per condition; one model size (9B) with LoRA and 50 steps; simulated environments with short episodes
(about 5--7 tool calls) and a fully described audit; heuristic labels (rule mentions, flagged tests) audited on
samples; judge-based misalignment scores with 200 agentic samples per model, enough to detect only large changes;
and a single reviewer, not blind to the model, for the action-level labels.

\paragraph{Future work.}
The behavior we could not elicit, escalation inside the environment, came in the incident from ingredients our
setup lacks: frontier-scale models, open-ended environments with unknown affordances, long horizons and persistent
state, uncertainty about how the grader works, many interacting agents, and training for persistence. Our session
environment shows that a write tool and an editable log are not enough at 9B. The natural next steps are larger
models on the same environments, rules generated per episode so that no single shortcut pays, grader uncertainty,
and additional seeds for the core results. \pending{SDF arms: whether giving the model a prior that links the
behavior to a broader category changes transfer.}

\paragraph{Author contributions.}
\todo{Yusuf: environments, RL runs, evaluations and analysis; John: Betley evaluation, SDF + RL on \ldots; Soham:
SDF + RL on \ldots}

\bibliography{references}
\bibliographystyle{icml2026}

\end{document}
